% REVISION_NOTE: Standalone revised contribution based on the earlier C23-429 attempt; the prior partial text is not overwritten.
% FULL_SOLUTION_SCOPE: {"target":"The specified finite-J, fixed-H synthetic-control interval construction, including pretreatment weight selection, validation informativeness, exact span, and contamination/deletion comparisons","result":"Simultaneous coverage theorem and explicit radius, nonparametric pretreatment certificate, excitation certificate, boundary counterexamples, exact-span and support selection rules","coverage":"All J>=1, T0>=0 for the universal bound, H>=1, alpha in (0,1), allowable rank and zero-noise boundary cases; additional hypotheses are stated explicitly only for refinements","dependencies":"Elementary Chernoff, union-bound, compactness and Gaussian integration arguments are supplied; background from Abadie (2021) and Cattaneo--Feng--Titiunik (2021)","human_verification":"None reported","unverified":"No independent human checking or formal verification; no claim of minimax-optimal length or novelty","novelty":"Primary literature/context checked; novelty of individual inequalities has not been established"}
\documentclass[11pt,a4paper]{article}
\usepackage[T1]{fontenc}
\usepackage[utf8]{inputenc}
\usepackage[margin=27mm]{geometry}
\usepackage{lmodern,microtype}
\usepackage{amsmath,amssymb,amsthm,mathtools}
\usepackage[hidelinks]{hyperref}
\setlength{\parindent}{0pt}
\setlength{\parskip}{5pt}
\newtheorem{theorem}{Theorem}[section]
\newtheorem{proposition}[theorem]{Proposition}
\newtheorem{lemma}[theorem]{Lemma}
\newcommand{\R}{\mathbb R}
\newcommand{\E}{\mathbb E}
\newcommand{\Prb}{\mathbb P}
\newcommand{\W}{\mathcal W}
\newcommand{\norm}[1]{\lVert#1\rVert}
\DeclareMathOperator{\supp}{supp}
\DeclareMathOperator*{\argmin}{arg\,min}
\title{Finite-sample coverage certificates for sparse synthetic controls}
\author{OpenAI GPT-6}
\date{11 October 2026}
\begin{document}
\maketitle
\textbf{Problem ID:} \texttt{C23-429} (JEL \texttt{C23}).
\textbf{Status:} Solved for the stated finite-donor construction and comparison target.
The results provide uniform coverage and explicit conditions for informative certificates; they do not assert that the proposed intervals are shortest or minimax optimal.

\section{A uniform interval that permits estimated donor weights}
The target is one treated region, $J\ge1$ donors, $T_0\ge0$ pretreatment periods, and $H\ge1$ post-treatment periods. Let $\W=\{w\in\R^J:w_j\ge0,\ \sum_j w_j=1\}$. All weights are selected from pretreatment observations (and possibly auxiliary randomization). The untreated outcomes obey
\[
 Y_{jt}(0)=\lambda_j^\top f_t+\eta_{jt}+\xi_{jt},\qquad
 r\le r_{\max},\quad\norm{\lambda_j}_2\le\Lambda,\quad\norm{f_t}_2\le F,
\]
where $|\eta_{jt}|\le a_{jt}$, and the mean-zero noises are mutually independent across regions and dates with
$\E e^{u\xi_{jt}}\le e^{u^2\sigma^2/2}$ for all real $u$.
The loadings and factors are deterministic model primitives (equivalently, the noise bound holds conditionally on fixed primitives). All regions are untreated before policy; afterward the observed donors satisfy $Y_{jt}=Y_{jt}(0)+s_{jt}$ with $|s_{jt}|\le b_{jt}$. The nonnegative envelopes $a_{jt},b_{jt}$ and constants $\Lambda,F,\sigma$ are declared in advance. Approximation errors and spillovers may otherwise be dependent on the data and noises.

For a post-treatment average, write a bar, and put
\[
 \begin{gathered}
 \Delta(w)=\lambda_0-\sum_{j=1}^Jw_j\lambda_j,\qquad
 Z_j=H^{-1}\sum_{t>T_0}\xi_{jt},\\
 A_+(w)=\bar a_0+\sum_jw_j\bar a_j,\qquad
 S_+(w)=\sum_jw_j\bar b_j,\\
 \tau=\bar Y_0(1)-\bar Y_0(0),\qquad
 \widehat\tau(w)=\bar Y_0(1)-\sum_jw_j\bar Y_j.
 \end{gathered}
\]
The exact decomposition is
\begin{equation}\label{eq:decomp}
 \widehat\tau(w)-\tau
 =\Delta(w)^\top\bar f
 +\bar\eta_0-\sum_jw_j(\bar\eta_j+\bar s_j)
 +Z_0-\sum_jw_jZ_j.
\end{equation}
The treated outcome under treatment cancels, so no restriction on its dependence with untreated noise is necessary. Coverage concerns the \emph{realized} average effect, including the treated unit's post-treatment counterfactual noise.

\begin{theorem}[Simultaneous finite-sample coverage]\label{thm:coverage}
Fix $0\le\delta_0<1$ and $0<\delta_1<1$ with $\delta_0+\delta_1\le\alpha<1$. Suppose a pretreatment-computable loading certificate $\widehat d(w)$ satisfies, uniformly over the declared model class,
\begin{equation}\label{eq:d-cert}
 \Prb\bigl\{\norm{\Delta(w)}_2\le\widehat d(w)
            \text{ for every }w\in\W\bigr\}\ge1-\delta_0.
\end{equation}
Define
\begin{equation}\label{eq:rad}
 R(w)=F\widehat d(w)+A_+(w)+S_+(w)
            +2\sigma\sqrt{\frac{\log(2J/\delta_1)}H}.
\end{equation}
Then
\begin{equation}\label{eq:simultaneous}
 \Prb\left\{\tau\in[\widehat\tau(w)-R(w),\widehat\tau(w)+R(w)]
                   \text{ for all }w\in\W\right\}
 \ge1-\delta_0-\delta_1.
\end{equation}
In particular, evaluating the interval at \emph{any} pretreatment-selected $\widehat w$ yields coverage at least $1-\alpha$ uniformly over the model class.
\end{theorem}
\begin{proof}
Independence and the sub-Gaussian moment bound give
$\E e^{u(Z_0-Z_j)}\le e^{u^2\sigma^2/H}$. Thus
\[
 \Prb\{|Z_0-Z_j|>x\}\le2\exp\{-Hx^2/(4\sigma^2)\},
\]
with the zero-noise case interpreted exactly. A union bound gives, with probability at least $1-\delta_1$,
$\max_j|Z_0-Z_j|\le2\sigma\sqrt{\log(2J/\delta_1)/H}$. On this \emph{same} event for every $w\in\W$,
\[
 \left|Z_0-\sum_jw_jZ_j\right|
 =\left|\sum_jw_j(Z_0-Z_j)\right|
 \le\sum_jw_j|Z_0-Z_j|
 \le2\sigma\sqrt{\frac{\log(2J/\delta_1)}H}.
\]
Intersect with the event in \eqref{eq:d-cert}. The factor term in \eqref{eq:decomp} is at most $F\widehat d(w)$, and the other terms have the respective pathwise envelopes. A union bound on these two \emph{events}, without assuming that they are independent, proves \eqref{eq:simultaneous}.
\end{proof}

The universal loading certificate is $\widehat d(w)=2\Lambda$ by the triangle inequality. Inserting $\delta_0=0$, $\delta_1=\alpha$ gives the explicit answer requiring \emph{no additional assumptions}:
\begin{equation}\label{eq:universal}
 \boxed{\quad
 C(\widehat w)=\widehat\tau(\widehat w)\ \pm\
 \left[2\Lambda F+\bar a_0+\sum_j\widehat w_j(\bar a_j+\bar b_j)
 +2\sigma\sqrt{\frac{\log(2J/\alpha)}H}\right].\quad}
\end{equation}
If $r_{\max}=0$, the factor-mismatch term vanishes. For comparison, the donor-free interval
\[
 \bar Y_0(1)\ \pm\ \left[\Lambda F+\bar a_0+
 \sigma\sqrt{\frac{2\log(2/\alpha)}H}\right]
\]
is also uniformly valid. Choosing the shorter of separately calibrated intervals requires joint error control; individual validity alone does not justify adaptive selection.

\section{Loading-mismatch certificates from pretreatment observations}
The universal radius is generally conservative. There are two ways to improve it while preserving \eqref{eq:d-cert} for \emph{all} candidate weights.

\textbf{Certificate without excitation.}
Suppose $T_0\ge1$ and $\delta_0>0$. Put
\[
 u_0=\sigma\sqrt{2\log\!\left(\frac{2(J+1)T_0}{\delta_0}\right)}.
\]
Pad all possible ranks to $r_{\max}$ coordinates. Define the \emph{observable} feasible set $\mathcal K_0$ of pretreatment loadings and factors by
\[
 \norm{\lambda_j}_2\le\Lambda,\quad \norm{f_t}_2\le F,
 \quad |Y_{jt}-\lambda_j^\top f_t|\le a_{jt}+u_0
 \quad (0\le j\le J,\ 1\le t\le T_0).
\]
For $\mathcal K_0\ne\varnothing$ let
\begin{equation}\label{eq:profile}
 \widehat d_{\mathcal K}(w)
 =\max_{(\lambda,f)\in\mathcal K_0}
          \left\|\lambda_0-\sum_jw_j\lambda_j\right\|_2,
\end{equation}
otherwise set $\widehat d_{\mathcal K}=2\Lambda$ and record that the calibrated feasibility set was empty. The set is compact, so the displayed maximum exists. The true factorization lies in $\mathcal K_0$ whenever all pretreatment noises satisfy $|\xi_{jt}|\le u_0$, an event of probability at least $1-\delta_0$ by a union bound. Therefore \eqref{eq:profile} satisfies \eqref{eq:d-cert}, without known factors, known exact rank, or factor excitation. Computing a \emph{global} maximum (or a certified upper bound) is essential: a local maximum is not a coverage certificate. Empty feasibility does not itself refute the stochastic model.

\textbf{A sharper certificate with factor excitation.}
Let $V$ be a prespecified pretreatment fitting or validation block of size $m\ge1$. Add the substantive assumption
\begin{equation}\label{eq:excite}
 m^{-1}\sum_{t\in V}f_tf_t^\top\succeq\kappa I_r,
 \qquad\kappa>0\text{ known},
\end{equation}
on the \emph{actual} factor dimension (not on artificial padded coordinates). Put
\[
 \begin{split}
 e_t(w)&=Y_{0t}-\sum_jw_jY_{jt},\\
 R_V(w)&=\left[m^{-1}\sum_{t\in V}e_t(w)^2\right]^{1/2},\\
 A_V(w)&=\left[m^{-1}\sum_{t\in V}
                 \left(a_{0t}+\sum_jw_ja_{jt}\right)^2\right]^{1/2}.
 \end{split}
\]
For $d=\min(r_{\max},m)>0$, define
\[
 K(d,x)=d+2\sqrt{dx}+2x,\qquad
 q_V=\sigma\sqrt{\frac{2K(d,\log(J/\delta_0))}m}.
\]

\begin{lemma}[Projected sub-Gaussian noise]\label{lem:proj}
If a random vector $X\in\R^n$ obeys
$\E e^{u^\top X}\le e^{s^2\norm u_2^2/2}$ for every $u\in\R^n$, and $P$ is a deterministic orthogonal projector of rank $d$, then
\[
 \Prb\{\norm{PX}_2^2>s^2K(d,x)\}\le e^{-x},\qquad x>0.
\]
\end{lemma}
\begin{proof}
Let $G\sim N(0,I_n)$ be independent of $X$. Gaussian integration and the assumed moment-generating bound give, for $0<t<1/(2s^2)$,
\[
 \E e^{t\norm{PX}_2^2}
 =\E_G\E_X e^{\sqrt{2t}(PG)^\top X}
 \le\E_Ge^{s^2t\norm{PG}_2^2}
 =(1-2s^2t)^{-d/2}.
\]
Consequently, for $U=\norm{PX}_2^2/s^2$ and $0<t<1/2$,
\[
 \log\E e^{t(U-d)}
 \le-dt-\frac d2\log(1-2t)
 \le\frac{dt^2}{1-2t}.
\]
Chernoff's inequality with $t=\sqrt{x}/(\sqrt d+2\sqrt{x})$ yields
$\Prb\{U>d+2\sqrt{dx}+2x\}\le e^{-x}$. The cases $s=0$ and $d=0$ are immediate.
\end{proof}

\begin{proposition}[Simultaneous excitation certificate]\label{prop:excite}
Under \eqref{eq:excite}, with probability at least $1-\delta_0$, simultaneously for every $w\in\W$,
\begin{equation}\label{eq:d-excite}
 \norm{\Delta(w)}_2\le
 \widehat d_V(w):=
 \min\left\{2\Lambda,\ \frac{R_V(w)+A_V(w)+q_V}{\sqrt\kappa}\right\}.
\end{equation}
If $r_{\max}=0$, the certificate is zero without needing \eqref{eq:excite}.
\end{proposition}
\begin{proof}
Let $B$ be the $m\times r$ matrix with rows $f_t^\top$, and let $P_V$ project onto its column space, of dimension at most $d$. Each vector $\xi_{0,V}-\xi_{j,V}$ has sub-Gaussian variance proxy $2\sigma^2I_m$. Lemma~\ref{lem:proj}, followed by a union bound over donors, gives
\[
 \max_j\frac{\norm{P_V(\xi_{0,V}-\xi_{j,V})}_2}{\sqrt m}
 \le q_V
\]
with probability at least $1-\delta_0$. Every weighted contrast is a convex combination of these donor contrasts. The latent residual $g(w)=B\Delta(w)$ belongs to the projector's range, hence projection and the triangle inequality imply simultaneously
\[
 \frac{\norm{g(w)}_2}{\sqrt m}
 \le R_V(w)+A_V(w)+q_V.
\]
Excitation gives $\norm{g(w)}_2/\sqrt m\ge\sqrt\kappa\norm{\Delta(w)}_2$. Combine the two inequalities with $\norm{\Delta(w)}_2\le2\Lambda$.
\end{proof}

The projection $P_V$ is used only to prove the bound; it is not estimated to evaluate \eqref{eq:d-excite}. Its stochastic allowance is of order
$\sigma\sqrt{(r_{\max}+\log(J/\delta_0))/m}$.
The actual fitting residual $R_V$ may still reflect noise. As the guarantee is simultaneous, $V$ may be used in the selection of the final weights. Multiple data-selected validation blocks require a joint event or an error-budget adjustment. An empirically guessed value of $\kappa$ does not replace the structural assumption \eqref{eq:excite}.

\section{Precisely when fitting and validation improve the interval}
Given a desired half-width $\varepsilon>0$, a sufficient, observable condition for a radius at most $\varepsilon$ is
\begin{equation}\label{eq:informative}
 F\widehat d_V(\widehat w)+A_+(\widehat w)+S_+(\widehat w)
 +2\sigma\sqrt{\frac{\log(2J/\delta_1)}H}\le\varepsilon.
\end{equation}
With $F>0$, the excitation-based loading allowance is strictly smaller than $2\Lambda F$ whenever
\[
 R_V(\widehat w)+A_V(\widehat w)+q_V
 <2\Lambda\sqrt\kappa.
\]
Thus informative fit requires excitation in the directions needed for extrapolation, sufficiently small residual and approximation envelopes, and non-dominant spillover and post-treatment-noise allowances. Enlarging a period in which the factors are unexcited does not accomplish this.

Full excitation is sufficient but not necessary for an average effect. Another valid \emph{additional} restriction is a known average-factor span:
\begin{equation}\label{eq:knownspan}
 \bar f=\sum_{t\in V}q_tf_t+v_f,
 \qquad\norm{v_f}_2\le\rho,
\end{equation}
where $q=(q_t)_{t\in V}$ is prespecified and deterministic. Form the validation-adjusted center
$\widehat\tau_q(w)=\widehat\tau(w)-\sum_{t\in V}q_te_t(w)$.
Substitution in \eqref{eq:decomp}, and the same donor-uniform noise argument, give a simultaneous interval of radius
\begin{align}\label{eq:span-radius}
 R_q(w)={}&\rho\widehat d(w)+A_+(w)+S_+(w)
 +\sum_{t\in V}|q_t|\left(a_{0t}+\sum_jw_ja_{jt}\right)\notag\\
 &+2\sigma\sqrt{\left(H^{-1}+\norm q_2^2\right)
                         \log(2J/\delta_1)}.
\end{align}
To verify it, replace the factor term $\Delta(w)^\top\bar f$ by
$\Delta(w)^\top v_f$, and note that each region's combined post-minus-pre noise $Z_j-\sum_{t\in V}q_t\xi_{jt}$ has variance proxy $\sigma^2(H^{-1}+\norm q_2^2)$, by independence across dates. The certificate event and this simultaneous noise event yield coverage $1-\delta_0-\delta_1$. If $\rho=0$, the loading certificate is unnecessary, and one may take $\delta_0=0$.

\section{Exact span, estimated weights, and the limits of pretreatment fit}
The exact-span variant needs care about what is assumed to balance. If the \emph{selected} weights themselves obey
$\lambda_0=\sum_j\widehat w_j\lambda_j$ almost surely, and spillovers are absent, then $\widehat d=0$, $S_+=0$ and Theorem~\ref{thm:coverage} gives
\[
 C(\widehat w)=\widehat\tau(\widehat w)\ \pm\
 \left[\bar a_0+\sum_j\widehat w_j\bar a_j
       +2\sigma\sqrt{\frac{\log(2J/\alpha)}H}\right].
\]
By contrast, the existence of $w^\star\in\W$ satisfying exact balance does not imply that an estimated $\widehat w$ equals $w^\star$. The selected mismatch obeys
\begin{equation}\label{eq:estimated}
 \norm{\Delta(\widehat w)}_2
 =\left\|\sum_j(w_j^\star-\widehat w_j)\lambda_j\right\|_2
 \le\Lambda\norm{\widehat w-w^\star}_1,
\end{equation}
but the unknown right side needs a certified bound. One direct implementation is to impose existence of a balancing vector $v\in\W$,
$\lambda_0=\sum_jv_j\lambda_j$, inside the feasible set $\mathcal K_0$ of \eqref{eq:profile}, and then maximize the mismatch of the \emph{actual chosen} weights. Proposition~\ref{prop:excite} is an alternative under excitation. Neither method treats the fitted weights as oracle weights.

\begin{proposition}[Perfect pretreatment fit need not imply small uncertainty]\label{prop:gap}
If $r_{\max}\ge1$ and $\Lambda F>0$, the model includes an observationally indistinguishable noiseless family with perfect pretreatment fit for any $T_0$, but with possible effects filling $[-\Lambda F,\Lambda F]$.
\end{proposition}
\begin{proof}
Take one donor with $\lambda_1=0$, a scalar factor $f_t=0$ on every pretreatment date and $f_t=F$ afterward, and set all noises, approximations, and spillovers to zero. Make the observed post-treatment treated outcomes zero. Allow $\lambda_0\in[-\Lambda,\Lambda]$. All observed data stay zero; the pretreatment donor residual vanishes identically; yet $\tau=-\lambda_0F$ fills the asserted interval.
\end{proof}

For deterministic interval rules and $1-\alpha>0$, coverage for all the indistinguishable models forces an interval to contain the entire set. With independent auxiliary randomization, endpoint coverage implies
\[
 \Prb\{\operatorname{length}(C)\ge2\Lambda F\}\ge1-2\alpha
 \quad(\alpha<\tfrac12),
\]
and integration of coverage over that compatible interval gives
$\E\operatorname{length}(C)\ge2\Lambda F(1-\alpha)$.
Even \emph{existence} of exact balance does not eliminate this example's logic: use two donors with loadings $-\Lambda$ and $+\Lambda$; any treated loading in between is exactly representable, but all pretreatment observations still vanish and the effect remains undetermined.

There is also an unavoidable prediction-noise scale at fixed $H$. Set all loadings, approximation errors, and spillovers to zero and take independent $N(0,\sigma^2)$ untreated noises, with the treated observed post-treatment outcomes fixed at zero. Then $\tau=-Z_0\sim N(0,\sigma^2/H)$ independently of the observed data. An interval of random length $L$ has conditional coverage at most
$g(L)=2\Phi(L\sqrt H/(2\sigma))-1$. Since $g$ is concave on $[0,\infty)$, Jensen's inequality yields
\begin{equation}\label{eq:noise-lower}
 \E L\ge\frac{2\sigma}{\sqrt H}\Phi^{-1}(1-\alpha/2).
\end{equation}
When $\sigma=0$ the lower bound is zero. Increasing $T_0$ alone does not remove either this uncertainty or the factor-direction obstruction.

\section{Sparse selection and the donor-contamination tradeoff}
It can be useful to make the noise radius sensitive to diversification. Besides the constant noise term in \eqref{eq:rad}, the following alternative is also simultaneous:
\begin{equation}\label{eq:ellipsoid}
 N_E(w)=\frac{\sigma}{\sqrt H}
 \sqrt{(1+\norm w_2^2)K(J,\log(1/\delta_1))}.
\end{equation}
Indeed, with $Z=(Z_0,\ldots,Z_J)$ and $P$ projecting onto the $J$-dimensional subspace perpendicular to $(1,\ldots,1)$, the contrast vector $(1,-w)$ lies in that subspace. Lemma~\ref{lem:proj} bounds $\norm{PZ}_2$, and Cauchy--Schwarz gives \eqref{eq:ellipsoid} simultaneously for every $w$. Either noise bound may be chosen in advance; taking their pointwise minimum requires joint calibration, such as a split of $\delta_1$.

For a prespecified available set $D$ and support limit $1\le k\le|D|$, a concrete sparse rule is to choose from pretreatment information
\begin{equation}\label{eq:sparselaw}
 \widehat w\in\argmin_{\substack{w\in\W,\ \supp(w)\subseteq D\\
                                    \norm w_0\le k}}
 \left\{F\widehat d(w)+\sum_jw_j(\bar a_j+\bar b_j)+N(w)\right\},
\end{equation}
with $N=q_+$ or $N_E$, and a deterministic tie-break. For the certificates above, the feasible set is a finite union of compact simplex faces and the objective is continuous, so a minimizer exists. For fixed small $J$, supports can be enumerated. Any feasible approximate minimizer also retains coverage; accuracy affects width, not the theorem. An $\ell_1$ penalty is constant on the simplex, $\norm w_1=1$, so it does not itself encourage sparse weights.

When exact balance exists in $\R^r$, there is a balancing representation using at most $r+1$ donors: if a representation has more than $r+1$ positive weights, the corresponding augmented vectors $(1,\lambda_j)$ are linearly dependent. A nonzero perturbation along that dependence preserves both the weight sum and the balanced loading. Move until one weight hits zero and iterate. This is an existence statement; it does not guarantee that noisy fitting attains that representation or that its contamination allowance is small.

For any original weights $w$ and refitted weights $v$, write
$R(w)=F\widehat d(w)+\bar a_0+\sum_jw_j\bar a_j+S_+(w)+N(w)$. Exactly,
\begin{align}\label{eq:deletion}
 R(v)-R(w)={}&F\{\widehat d(v)-\widehat d(w)\}
 +\sum_j(v_j-w_j)\bar a_j\notag\\
 &+\{S_+(v)-S_+(w)\}+\{N(v)-N(w)\}.
\end{align}
If the remaining donors used by $v$ are uncontaminated, deletion reduces the certified half-width \emph{if and only if}
\begin{equation}\label{eq:delete-condition}
 S_+(w)>F\{\widehat d(v)-\widehat d(w)\}
 +\sum_j(v_j-w_j)\bar a_j+N(v)-N(w).
\end{equation}
This states the spillover protection versus the loss of balance, approximation control, and noise diversification. For example, if $\lambda_0=\lambda_1=0$, $\lambda_2=\delta$ in a scalar model, donor~1 has spillover envelope $b$, and donor~2 is clean, the deterministic single-donor allowances are $b$ and $F|\delta|$. Their noise allowances agree, so switching to the clean donor reduces the radius precisely when $b>F|\delta|$.

For the \emph{same} globally defined radius and bounds, restricting the donor set cannot decrease the optimized minimum:
\[
 \inf_{\substack{w\in\W,\ \supp(w)\subseteq D}}R(w)
 \ \ge\ \inf_{w\in\W}R(w).
\]
A contamination-aware optimizer over the full pool can already assign zero weight to a contaminated donor. Removal can improve on a fit-only rule or hedge against larger uncertainty about contamination, but not beat the globally minimized common criterion. If donors are deleted adaptively, one cannot simply replace $J$ in the simultaneous bound by the selected subset size without recalibration.

\section{Completion Estimate}
\noindent\textbf{COMPLETION: 100\%}
The finite-donor problem as stated is resolved constructively: \eqref{eq:universal} has uniform coverage without additional assumptions, \eqref{eq:profile} and \eqref{eq:d-excite} permit valid data-dependent reductions, \eqref{eq:informative}--\eqref{eq:span-radius} give explicit informative-validation criteria, and \eqref{eq:deletion}--\eqref{eq:delete-condition} settle the tradeoff for the proposed certified radii. The exact-span variants and counterexamples are covered as well. The completion estimate refers \emph{only} to the original specified target, not to dependent-shock extensions, necessary-and-sufficient minimax informativeness conditions, or shortest possible intervals. It is an LLM claim, not an independently verified mathematical verdict.

\section*{Sources and verification}
The mathematical arguments above were generated and checked by OpenAI GPT-6 on 11 October 2026, using the user-provided problem and earlier attempt. LaTeX was compiled to check syntax and references. No independent human mathematical proof check or formal verification is reported; source checking does not establish novelty. The original problem is an editorial finite-donor adaptation of a methodological agenda and should not be read as a theorem already stated by its source. The prior attempt's stronger independence assumption was removed by using simultaneous noise events; no assumption that pretreatment errors are independent of future noises is required. No empirical effect is claimed.

\begin{thebibliography}{9}
\bibitem{abadie}
Alberto Abadie (2021), \emph{Using Synthetic Controls: Feasibility, Data Requirements, and Methodological Aspects},
\emph{Journal of Economic Literature} 59(2), 391--425.
\url{https://doi.org/10.1257/jel.20191450}; primary text:
\url{https://economics.mit.edu/sites/default/files/publications/jel.20191450.pdf}.
The methodological and inference discussion motivates the stated finite-donor adaptation.

\bibitem{cft}
Matias D. Cattaneo, Yingjie Feng, and Rocio Titiunik (2021),
\emph{Prediction Intervals for Synthetic Control Methods},
\emph{Journal of the American Statistical Association} 116(536), 1865--1880.
\url{https://doi.org/10.1080/01621459.2021.1979561};
\url{https://mdcattaneo.github.io/papers/Cattaneo-Feng-Titiunik_2021_JASA.pdf}.
This is existing synthetic-control inference literature, not a source for the displayed finite-donor certificates.
\end{thebibliography}
\end{document}
